\section{Введение}

\subsection{Актуальность}

Ленты коротких видео стали одним из основных форматов потребления контента: сервисы ежедневно обслуживают миллионы пользователей, а количество зарегистрированных взаимодействий (просмотров, положительных отметок, пересылок) измеряется миллиардами событий.
В таких условиях качество ленты напрямую определяется системой ранжирования --- моделью, которая для каждого пользователя упорядочивает кандидатов-видео по ожидаемой полезности.

Построение промышленной системы ранжирования требует решения двух связанных задач.
Во-первых, это задача инженерии данных (Data Engineering): исходные журналы взаимодействий имеют объем в десятки миллионов строк, требуют очистки, агрегации и трансформации в признаковое описание, пригодное для машинного обучения, причем обработка должна выполняться в потоковом режиме, не загружая весь объем в память.
Во-вторых, это задача машинного обучения: выбор целевой переменной, отражающей интерес пользователя, построение признаков без утечек будущего, обучение ранжирующей модели и ее корректная оценка офлайн-метриками ранжирования.

Отдельное требование современной ML-разработки --- воспроизводимость экспериментов.
Каждая стадия обработки данных и каждое обучение модели должны фиксироваться в системе отслеживания экспериментов вместе с конфигурацией, метриками и артефактами.
В работе для этого используется платформа ClearML~\cite{clearml-docs}, развернутая локально в Docker.

\subsection{Постановка задачи}

В рамках производственной практики были поставлены следующие задачи:

\begin{enumerate}
    \item Провести анализ требований к обработке и анализу больших данных для задачи ранжирования коротких видео.
    \item Разработать пайплайн сбора, хранения и предобработки данных с использованием инструментов Data Engineering (Polars, PyArrow).
    \item Настроить процессы очистки, агрегации и трансформации данных для последующего анализа и применения ML-моделей.
    \item Провести анализ данных, рассчитать ключевые метрики, сформировать признаки и подготовить датасеты для машинного обучения.
    \item Обучить и оценить ранжирующую модель с использованием ClearML, CatBoost~\cite{catboost-docs} и PyTorch~\cite{pytorch-docs}; обучение итоговой модели выполнить на GPU.
    \item Протестировать реализованные пайплайны и ML-модели автоматизированными тестами.
    \item Подготовить итоговый отчет с описанием архитектуры решения, выбранного стека и результатов.
\end{enumerate}

\subsection{Задачи}

\begin{enumerate}
    \item Изучить структуру исходного датасета взаимодействий пользователей с короткими видео: недельные файлы событий, метаданные пользователей и видео, контентные эмбеддинги.
    \item Реализовать этап индексации данных с сортировкой недельных файлов по глобальному временнóму порядку.
    \item Реализовать потоковую очистку событий и построение градуированной целевой переменной (релевантности) из неявной обратной связи.
    \item Спроектировать и вычислить три группы признаков --- признаки пользователя, признаки видео и парные (пользователь--видео), включая признак близости контентного эмбеддинга видео к эмбеддинг-профилю пользователя (расчет на PyTorch).
    \item Выполнить временнóе разбиение данных на обучающую, валидационную и тестовую выборки и собрать датасеты в формате, требуемом ранжирующей моделью CatBoost.
    \item Реализовать собственный модуль офлайн-метрик ранжирования (NDCG@k, MAP@k, MRR@k, HitRate@k, Recall@k, Precision@k, GAUC) в векторизованном виде.
    \item Обучить модель \texttt{CatBoostRanker} (функция потерь YetiRank~\cite{yetirank}) с расчетом метрик на валидации в процессе обучения и логированием в ClearML.
    \item Покрыть пайплайн и модель автоматизированными тестами на синтетических данных.
\end{enumerate}

Исходный код решения вместе с используемыми данными опубликован в открытом репозитории:
\url{https://github.com/seregamikhailov/ml-training-catboost}.

\subsection{Сценарии использования (Use Cases)}

UML-диаграмма прецедентов разработанного решения приведена на рисунке~\ref{fig:usecase}.
Актерами выступают инженер (запускает пайплайн и обучение), исследователь (анализирует результаты) и сервер ClearML как внешняя система отслеживания экспериментов.

\begin{figure}[H]
\centering
\begin{tikzpicture}
  % --- граница системы
  \draw (0,0) rectangle (9.6,9.4);
  \node[font=\small\bfseries, anchor=north] at (4.8,9.2)
    {ML-пайплайн ранжирования коротких видео};

  % --- прецеденты
  \node[usecase] (uc2) at (4.8,7.6) {UC-2. Проверочный прогон\\на подвыборке};
  \node[usecase] (uc1) at (4.8,5.5) {UC-1. Полный прогон\\пайплайна};
  \node[usecase] (uc3) at (4.8,3.4) {UC-3. Обучение итоговой\\модели на GPU};
  \node[usecase] (uc4) at (4.8,1.3) {UC-4. Анализ\\эксперимента};

  % --- актеры
  \pic at (-2.2,7.2) {stickman};
  \node[actorname] at (-2.2,5.3) {Инженер};
  \pic at (-2.2,2.4) {stickman};
  \node[actorname] at (-2.2,0.5) {Исследователь};
  \pic at (12.0,5.2) {stickman};
  \node[actorname] at (12.0,3.3) {Сервер\\ClearML};

  % --- ассоциации актеров с прецедентами
  \draw[assoc] (-1.7,6.5) -- (uc2.west);
  \draw[assoc] (-1.7,6.5) -- (uc1.west);
  \draw[assoc] (-1.7,6.5) -- (uc3.west);
  \draw[assoc] (-1.7,1.7) -- (uc4.west);
  \draw[assoc] (11.5,4.5) -- (uc1.east);
  \draw[assoc] (11.5,4.5) -- (uc3.east);
  \draw[assoc] (11.5,4.5) -- (uc4.east);

  % --- отношения между прецедентами
  % проверочный прогон — вариация полного (extend),
  % полный прогон включает обучение модели (include)
  \draw[ucdep] (uc2.south) -- node[right, font=\scriptsize] {$\ll$extend$\gg$} (uc1.north);
  \draw[ucdep] (uc1.south) -- node[right, font=\scriptsize] {$\ll$include$\gg$} (uc3.north);
\end{tikzpicture}
\caption{UML-диаграмма прецедентов разработанного решения}
\label{fig:usecase}
\end{figure}

\textbf{UC-1. Полный прогон пайплайна.}
Инженер запускает команду \texttt{python -m src.pipeline}.
Система последовательно выполняет пять этапов: индексацию данных, очистку, построение признаков, сборку датасетов и обучение модели.
Каждый этап создает отдельный эксперимент в ClearML с конфигурацией, статистиками и артефактами.

\textbf{UC-2. Быстрый проверочный прогон.}
Инженер ограничивает объем данных параметрами конфигурации (\texttt{max\_shards}, \texttt{user\_fraction}, \texttt{max\_rows}) и прогоняет пайплайн на подвыборке за минуты, проверяя работоспособность всех этапов перед полномасштабным запуском.

\textbf{UC-3. Обучение итоговой модели на GPU.}
Инженер запускает этап \texttt{train} на сервере с GPU по полному конфигу (\texttt{task\_type: GPU}).
На машине без CUDA обучение автоматически продолжается на CPU (механизм \texttt{allow\_cpu\_fallback}), что позволяет использовать один и тот же код локально и на сервере.

\textbf{UC-4. Анализ эксперимента.}
Исследователь открывает веб-интерфейс ClearML, сравнивает кривые метрик по итерациям обучения, изучает важности признаков и итоговые офлайн-метрики на валидации и тесте, скачивает артефакт модели.
